\section{Primitives}\label{sec:primitives}

Let \(A\) be the type of hidden presentations, with semantic interpretation
\(M_x\) for each \(x\in A\). Write \(x\simeq y\) for homotopy equivalence
of \(M_x\) and \(M_y\) in Helfer's sense \cite[Definition~7.1]{helfer2019fol}.
Assume that this equivalence preserves the chosen closed tests. Let \(P\) be
the type of public class data and choose a classifier \(\pi:A\to P\) that factors through their
closed theory. Then
\[
  x\simeq y\quad\Longrightarrow\quad\pi(x)=\pi(y).
\]
The typed data and the invariance premise are made precise in
\cref{def:typed-hcp-realization,prop:helfer-hcp-bridge}.
For a realized public value \(\alpha\), write
\[
  \alpha\in\image(\pi),\qquad
  A_\alpha=\pi^{-1}(\alpha)=\{x\in A:\pi(x)=\alpha\}.
\]
For \(x\in A\), we also write
\([x]_\pi:=A_{\pi(x)}\). Thus, whenever \(\alpha=\pi(x)\), the two
descriptions agree: \(A_\alpha=[x]_\pi\). We call the classifier
\emph{complete for the declared homotopy equivalence} when
\(\pi(x)=\pi(y)\) implies \(x\simeq y\) for all \(x,y\in A\).

We attach Hodge data to each presentation using finite graded cell carriers,
integer incidence numbers and face-to-coface linear stalk maps.

Fix a finite coordinate set \(N\) and let \(V=\mathbb R^N\) have its Euclidean
inner product. Each presentation has finite-dimensional real inner-product
cochain spaces \(C_x^0,C_x^1,C_x^2\). To compare presentations using one public
cochain vector, fix isometric coordinate charts
\(\iota_x:C_x^1\overset{\sim}{\longrightarrow}V\).
We identify degree-one vectors and maps with their transported versions in
\(V\); the charts are part of the comparison data.

The restriction maps \(R^x_{\sigma\to\tau}\) for \(\sigma\prec\tau\),
including vertex-to-edge and edge-to-face maps, determine adjacent
coboundaries
\[
  \delta_{0,x}:C_x^0\to V,
  \qquad
  \delta_{1,x}:V\to C_x^2,
  \qquad
  \delta_{1,x}\delta_{0,x}=0.
\]
In these coordinates, define the degree-one Hodge--Laplacian, its harmonic
subspace and its orthogonal projector by
\[
\begin{aligned}
  \Delta_x&=\delta_{0,x}\delta_{0,x}^\ast+
    \delta_{1,x}^\ast\delta_{1,x},\\
  \mathcal H_x&=\ker(\Delta_x),\qquad
    P_{\mathcal H_x}:V\to\mathcal H_x\subseteq V.
\end{aligned}
\]
The stars denote adjoints. When the operative representative is fixed, we
suppress \(x\) and also write \(\Delta_1\) or \(\Delta\) for its degree-one
Laplacian, and \(P_H\) for its harmonic projector.

The algebraic cellular model above adapts the finite Hodge setup of
\cite{hansenGhrist2018sheaves} without assuming a regular CW or simplicial
realization for every admitted incidence structure. Such a realization and
the separate simplicial-coherence conditions require additional data.

\begin{definition}[Matrix harmonicity]\label{def:matrix-harmonicity}
  For a hidden linear Hodge operator \(\Delta\), a cochain
  \(h\in\mathbb R^N\)
  is \(\Delta\)-harmonic when
  \[
    \Delta h=0.
  \]
  Thus the harmonic kernel is \(\ker(\Delta)\). In the projection model this
  agrees with the condition \(h\in W\) for the kernel-normalized operator
  \(\Delta_W=I-P_W\).
\end{definition}

\begin{definition}[Harmonic projection]\label{def:projection-laplacian}
  For a subspace \(W\le V\), let \(P_W:V\to W\subset V\) denote the
  orthogonal projection. Its associated kernel-normalized operator is
  \(\Delta_W=I-P_W\), with \(\ker\Delta_W=W\).
\end{definition}

\begin{proposition}[Hodge decomposition surface]\label{prop:hodge-decomposition}
  Let the three consecutive cochain spaces be finite-dimensional real inner
  product spaces, and let the displayed stars denote adjoints. Suppose
  \[
    d_{\mathrm{next}}d_{\mathrm{prev}}=0,
    \qquad
    \Delta=d_{\mathrm{prev}}d_{\mathrm{prev}}^\ast
      +d_{\mathrm{next}}^\ast d_{\mathrm{next}} .
  \]
  Then every cochain \(b\) admits a decomposition
  \[
    b=h+e_1+e_2
  \]
  with
  \[
    \Delta h=0,\qquad
    e_1\in\image(d_{\mathrm{prev}}),\qquad
    e_2\in\image(d_{\mathrm{next}}^\ast),
  \]
  and
  \[
    \langle h,e_1\rangle=\langle h,e_2\rangle=\langle e_1,e_2\rangle=0.
  \]
  The decomposition is unique.
\end{proposition}

\begin{proof}
  Each summand defining \(\Delta\) is self-adjoint, so \(\Delta\) is
  self-adjoint. In finite dimension,
  \(\image(\Delta)=(\ker\Delta)^\perp\), and hence the ambient cochain
  space decomposes as
  \[
    \ker(\Delta)\oplus\image(\Delta).
  \]
  Since
  \[
    \image(\Delta)
    \subseteq
    \image(d_{\mathrm{prev}})+\image(d_{\mathrm{next}}^\ast),
  \]
  every cochain \(b\) can be written as \(h+e_1+e_2\) with \(h\in\ker\Delta\),
  \(e_1\in\image(d_{\mathrm{prev}})\), and
  \(e_2\in\image(d_{\mathrm{next}}^\ast)\). From \(\Delta h=0\), one obtains
  \[
    \|d_{\mathrm{prev}}^\ast h\|^2+\|d_{\mathrm{next}}h\|^2=0,
  \]
  hence \(d_{\mathrm{prev}}^\ast h=0\) and \(d_{\mathrm{next}}h=0\). These
  identities give \(\langle h,e_1\rangle=0\) and
  \(\langle h,e_2\rangle=0\), while \(d_{\mathrm{next}}d_{\mathrm{prev}}=0\)
  gives \(\langle e_1,e_2\rangle=0\). Hence the harmonic, exact, and
  coexact subspaces are pairwise orthogonal. Their sum is therefore direct,
  which makes the three components unique.
\end{proof}

\subsection{Hodge Encryption and Correctness}

For a hidden representative \(x\), define its noise space by
\[
  \mathcal E_x=\mathcal H_x^\perp.
\]
The Hodge decomposition identifies the noise space as the orthogonal direct
sum
\[
  \mathcal E_x
  =
  \image(\delta_{0,x})
  \oplus^\perp
  \image(\delta_{1,x}^{\ast}),
\]
and gives the orthogonal decomposition
\[
  C_x^1=\mathcal H_x\oplus^\perp\mathcal E_x.
\]
Consequently,
\[
\begin{aligned}
  \operatorname{combine}_x:
    \mathcal H_x\times\mathcal E_x&\longrightarrow C_x^1,
  &(h,e)&\longmapsto h+e,\\
  \operatorname{split}_x:C_x^1&\longrightarrow
    \mathcal H_x\times\mathcal E_x,
  &b&\longmapsto
    \bigl(P_{\mathcal H_x}b,P_{\mathcal E_x}b\bigr)
\end{aligned}
\]
are inverse linear bijections. Orthogonality also gives
\[
  \|h+e\|^2=\|h\|^2+\|e\|^2.
\]

This bijection is the algebraic encryption/decryption construction. Given a
harmonic message \(h\in\mathcal H_x\), randomized Hodge encryption samples, or
otherwise chooses, admissible noise \(e\in\mathcal E_x\) and sets
\[
  \mathsf{Enc}_x(h;e)=
  \operatorname{combine}_x(h,e)=h+e.
\]
Full coordinate recovery is \(\operatorname{split}_x\); message decryption
retains its harmonic coordinate:
\[
  \mathsf{Dec}_x(b)
  =
  \operatorname{pr}_{\mathcal H_x}
    \bigl(\operatorname{split}_x(b)\bigr)
  =
  P_{\mathcal H_x}b.
\]
For every admissible noise value,
\[
  \mathsf{Dec}_x\bigl(\mathsf{Enc}_x(h;e)\bigr)=h,
  \qquad
  \operatorname{split}_x\bigl(\operatorname{combine}_x(h,e)\bigr)=(h,e).
\]
The bijection concerns the full message--noise pair \((h,e)\). Choosing a
noise distribution makes \(\mathsf{Enc}_x\) randomized, while
\(\mathsf{Dec}_x\) recovers the message for every admissible noise value.
Both maps depend on the operative representative \(x\).

In the \HE{} public-sample model, a presentation \(x\) produces a cochain
\[
  b=h+e,
  \qquad
  h\in\mathcal H_x,
  \qquad
  e\in\image(\delta_{0,x})\oplus\image(\delta_{1,x}^\ast).
\]
The message \(h\) is a harmonic representative of its degree-one cohomology
class, and \(e\) is Hodge-structured noise. When \(x\) is also equipped with
an explicit harmonic-lattice presentation, write
\(\Lambda_x\subseteq\mathcal H_x\) for its discrete lattice and choose
\(h\in\Lambda_x\). In this integral profile, the public fiber indexes a
family of possible embedded harmonic lattices; the operative hidden member
selects one of them.
